% ======================================================================
%  第 3 讲  扭化：为什么要把超对称「扭」掉
%  页锚以 B 版原文（\cite{kw06}）的印刷页为准
% ======================================================================
\section{扭化：为什么要把超对称「扭」掉}
\label{sec:twist}

\begin{goals}
\begin{enumerate}
\item 说清楚弯曲流形上原来的超荷为什么不够用，拓扑场论到底向超荷要什么。
\item 说清楚 GL 扭化选了哪个子群，为什么恰好是它把六个标量排成一个一形式加两个标量。
\item 说清楚参数 $t\in\PP^1$ 与典范参数 $\Psi$ 的关系，以及为什么只有 $\Psi$ 有物理意义。
\end{enumerate}
\end{goals}

\subsection{弯曲流形上的麻烦}

先回忆平坦 $\RR^4$ 上的情形。$\mathcal N=4$ 超对称 Yang--Mills 理论有 16 个超荷；在超对称场论里，能动张量可以写成
\begin{equation}
T_{\mu\nu}=\{Q,\,G_{\mu\nu}\}
\label{eq:T-exact}
\end{equation}
的形式（$G_{\mu\nu}$ 是某个张量算符，等号在适当的改善项意义下成立）。只要真空保持 $Q$，$T_{\mu\nu}$ 就是 $Q$-恰当的，基态能量为零，配分函数对度量的整体缩放不敏感。这是拓扑场论最朴素的来源：路径积分会「局部化」到 $Q$-不变构型上，而 $Q$-恰当的能动张量保证度量形变不改变关联函数。

麻烦在于，这个论证默认 $Q$ 是平坦空间里的整体守恒荷。把理论放到一般弯曲流形 $M$ 上，超荷不再自动是整体定义的标量：它们按 $M$ 的旋量丛变换，沿 $M$ 平行移动时会互相混合。一般度规的和乐群会把 16 个超荷搅在一起，通常一个也留不下来，更谈不上从中挑出一个满足 \eqref{eq:T-exact} 的 $Q$。换句话说，普通超对称在弯曲流形上\emph{自动}掉了，理论并不因此变成拓扑的。\src{B p.16}

要让理论「拓扑」，就必须换一个在任意 $M$ 上都存在的超荷。做法不是改物理，而是改我们称呼 Lorentz 群的方式：把 $M$ 切丛的结构群 $\mathrm{Spin}(4)$ 对角地嵌入 $\mathrm{Spin}(4)\times\mathrm{Spin}(6)$，其中 $\mathrm{Spin}(6)$ 是 $\mathcal N=4$ 理论的 $R$-对称群。这个操作叫\term{twisting}{扭化}。扭化之后，场的自旋按新的对角群来数，原来按 $\mathrm{Spin}(6)$ 指标变换的场被重新组合成 $M$ 上的普通张量场；更重要的是，16 个超荷里会出现若干在 $M$ 上整体定义的标量分量——它们才是候选的拓扑超荷。在平坦空间上，扭化前后的理论只是同一批场的重新打包，物理内容并没有变。

拓扑场论的判据可以写得很干净：存在一个标量超荷 $Q$，满足
\begin{equation}
Q^2=0,\qquad T_{\mu\nu}=\{Q,\,b_{\mu\nu}\},
\label{eq:top-criterion}
\end{equation}
于是一切度量变分都是 $Q$-恰当的，关联函数与度量无关。这个判据在 \cite{fr09} 里有明确表述（\src{D p.13}）。需要说明的是，\cite{kw06} 并没有从这个判据出发去证明什么；原文的做法是\emph{先构造}：挑一个扭化，验证 $Q^2=0$，再检查得到的理论确实有拓扑性。判据是事后的组织原则，不是原文的证明路线。这一点值得先记住，因为本讲后面几处「为什么恰好是这个扭化」都只有构造性的回答。

在 $\mathcal N=4$ 上，$R$-对称群是 $\mathrm{Spin}(6)\cong SU(4)$，比 $\mathcal N=2$ 的 $U(1)_R$ 大得多，因此可选的扭化不止一种。原文 p.17 说有\emph{三种}不等价的扭化：前两种（其中一种在文献 [45] 里讨论过）与 Donaldson 理论密切相关，给出瞬子型的拓扑不变量；第三种在 [52--57] 里被研究过，却一直没有找到用武之地，而它恰好是与几何 Langlands 有关的那一个，原文称之为 \term{GL twist}{GL 扭化}。Donaldson--Witten 的类比是这样来的：$\mathcal N=2$ 超对称 Yang--Mills 经过一次 $U(1)_R$ 扭化成为拓扑 Yang--Mills 理论，其路径积分局部化在瞬子上，关联函数给出四维流形的 Donaldson 不变量。第三种扭化可以看作这个故事的 $\mathcal N=4$ 版本，只不过它多留下了一整个参数族。

\subsection{GL 扭化：选哪个子群}

把 $\mathcal N=4$ 的场按 $\mathrm{Spin}(4)\times\mathrm{Spin}(6)$ 写出来：规范场 $A_\mu$ 是 $(\mathbf 2,\mathbf 2;\mathbf 1)$，六个标量 $\varphi_I$ 是 $(\mathbf 1,\mathbf 1;\mathbf 6)$，费米子落在 $(\mathbf 2,\mathbf 1;\mathbf 4)\oplus(\mathbf 1,\mathbf 2;\bar{\mathbf 4})$ 里。GL 扭化选取一个同态
\begin{equation}
\kappa:\ \mathrm{Spin}(4)\longrightarrow\mathrm{Spin}(6),
\label{eq:kappa}
\end{equation}
它把 $\mathrm{Spin}(4)$ 的旋量表示 $\mathbf 4$ 送到 $\mathrm{Spin}(6)$ 的旋量表示 $\mathbf 4$，然后令
\begin{equation}
\mathrm{Spin}'(4)=(1\times\kappa)\bigl(\mathrm{Spin}(4)\bigr)\subset \mathrm{Spin}(4)\times\mathrm{Spin}(6)
\label{eq:spinx}
\end{equation}
为新的 Lorentz 群。$\kappa$ 在 $\RR^4$ 上与原来的 $\mathrm{Spin}(4)$ 作用相同，但在规范场上作用不同——这正是「扭」这个字的全部含义。\src{B p.17}

扭化能不能给出拓扑超荷，判据是正手征旋量空间 $S_+$ 里有没有非零的不变向量。对 GL 扭化，答案是肯定的，而且不止一个：$S_+$ 里有两个线性无关的不变旋量，记作 $\varepsilon_\ell,\varepsilon_r$，分别来自 $(\mathbf 2,\mathbf 1;\mathbf 4)$ 与 $(\mathbf 1,\mathbf 2;\bar{\mathbf 4})$。\src{B p.18} 相应地，扭化后的理论保留两个标量超荷 $Q_\ell,Q_r$，任意组合
\begin{equation}
Q=u\,Q_\ell+v\,Q_r,\qquad (u,v)\neq(0,0)
\label{eq:Q-combo}
\end{equation}
都是标量、都在任意 $M$ 上有定义。为什么 $Q^2=0$？原文的论证是：$Q^2$ 作为算符必须与新的 Lorentz 群交换，而 \src{B p.18 eq (2.15)} 的右边不存在 $\mathrm{Spin}'(4)$-不变的算符，唯一可能是零；写成反对易子就是 $Q^2=\tfrac12\{Q,Q\}=0$。\src{B p.18} 这个论证只针对平移部分。把规范部分也算进来，有
\begin{equation}
\delta_T^2=-\ii\,(u^2+v^2)\,L_\sigma,
\label{eq:deltaT2}
\end{equation}
其中 $L_\sigma$ 是场依赖的规范变换（$\sigma$ 是后面要出现的标量）。因此严格平方为零只在 $u^2+v^2=0$ 时成立；对一般的 $t$，只能说 $Q^2$ 是「模掉一个规范变换」为零。\src{B p.31}\src{B p.33} 后文会看到，$u^2+v^2=0$ 恰是 $t=\pm\ii$ 那两个点。

为什么要选这个 $\kappa$，而不是前两种扭化？原文没有给出「非它不可」的证明，这是需要诚实说明的地方：它只构造并指出这一种与几何 Langlands 相关。不过构造本身已经透露了理由。按 $\mathrm{Spin}'(4)$ 分解，六个标量所在的空间 $\wedge^2\mathbf 4$ 拆成
\begin{equation}
\mathbf 6_v=\wedge^2\mathbf 4
=(\mathbf 2,\mathbf 2)_0\oplus(\mathbf 1,\mathbf 1)_{2}\oplus(\mathbf 1,\mathbf 1)_{-2},
\label{eq:6-decomp}
\end{equation}
下标的 $K$ 是 $\mathrm{Spin}(6)$ 里由 $\Gamma_{89}$ 生成的 $U(1)$ 荷。第一项 $(\mathbf 2,\mathbf 2)_0$ 正是一个一形式，后两项是两个标量。也就是说，扭化把「六个标量」自动重组为「一个伴随值一形式加两个复标量」；而一形式 $\varphi=\varphi_\mu\dd x^\mu$ 与规范场 $A_\mu$ 可以合并成复联络 $\mathcal A=A+\ii\varphi$。当 $t=\pm\ii$ 时，后面的 BPS 方程恰好是 $\mathcal A$ 平坦、外加 $D{*}\varphi=0$，这正是 Hitchin 方程；而平坦 $G_{\mathbb C}$ 联络与 Hitchin 模空间正是几何 Langlands 的舞台。前两种扭化给不出这样的「一形式加平坦复联络」结构，这就是 GL 扭化与几何 Langlands 相连的结构性原因。

\subsection{扭化后场长什么样}

把扭化后的场列成表，会看得更清楚。玻色子的内容如表~\ref{tab:bosons} 所示：
\begin{table}[htbp]
\centering
\caption{扭化后玻色场的 $\mathrm{Spin}'(4)$ 表示与 $K=\Gamma_{89}$ 荷。前两行同属 $(\mathbf 2,\mathbf 2)_0$，扭化后在 Lorentz 群下不可区分；区别在来源。}
\label{tab:bosons}
\begin{tabular}{@{}llcl@{}}
\toprule
场 & $\mathrm{Spin}'(4)$ 表示 & $K=\Gamma_{89}$ & 来源 \\
\midrule
$A_\mu$ & $(\mathbf 2,\mathbf 2)_0$ & $0$ & 四维规范场 \\
$\varphi_\mu$ & $(\mathbf 2,\mathbf 2)_0$ & $0$ & $\mathbf 6$ 中的 $\mathbf 4$ 个标量 \\
$\sigma$ & $(\mathbf 1,\mathbf 1)_2$ & $2$ & $\mathbf 6$ 中的单态 \\
$\bar\sigma$ & $(\mathbf 1,\mathbf 1)_{-2}$ & $-2$ & $\mathbf 6$ 中的单态 \\
\bottomrule
\end{tabular}
\end{table}
\src{B p.19} 两个一形式都是 $(\mathbf 2,\mathbf 2)_0$，在扭化后的 Lorentz 群下无法区分，它们的区别在来源：$A_\mu$ 是规范场，$\varphi_\mu$ 来自标量。$\sigma$ 与 $\bar\sigma$ 是复标量；它们的 $K$ 荷为 $\pm2$，这解释了为什么 $\sigma$ 在超荷下不变：超对称变分要把 $\sigma$ 变成 $K=3$ 的费米子，而场里根本没有 $K=3$ 的分量，于是 $\delta_T\sigma=0$。\src{B p.24}

费米子的 16 个分量按 $K$ 分块，只有 $K=-\ii$ 的两个分量与不变超荷对应，其余分量两两配对，在扭化后的理论里不产生额外的标量超荷。\src{B p.23} 一个容易忽略的细节是：两个不变超荷的 $K$ 荷\emph{相同}。这一点直接决定了参数族的形状。若两个超荷的 $K$ 荷不同，$U(1)_K$ 对称性会把其中一个（或某个线性组合）挑出来，$(u,v)$ 的整个 $\PP^1$ 就被压成一个点；正因为 $K$ 荷相同，两个超荷才平权，族才是完整的 $\PP^1$。\src{B p.22}

扭化后的超对称变换把 $A_\mu$ 与 $\varphi_\mu$ 混合，把 $\sigma$ 单独留下。写出来最能说明问题的是 $t=\pm\ii$ 时的 BPS 方程：
\begin{equation}
\bigl(F-\varphi\wedge\varphi+t\,D\varphi\bigr)^+=0,\qquad
\bigl(F-\varphi\wedge\varphi-t^{-1}D\varphi\bigr)^-=0,\qquad
D{*}\varphi=0,
\label{eq:bps}
\end{equation}
其中 $\pm$ 是四维流形上的自对偶/反自对偶投影。当 $t=\pm\ii$ 时，前两条方程等价于
\begin{equation}
F(\mathcal A)=0,\qquad \mathcal A=A+\ii\varphi,
\label{eq:flat}
\end{equation}
即\emph{平坦 $G_{\mathbb C}$ 联络}，再加上 $D{*}\varphi=0$。其解模规范变换，就得到平坦联络模空间的一部分（\src{B p.25}）：
\begin{equation}
Y_M(G)=\operatorname{Hom}\bigl(\pi_1(M),G_{\mathbb C}\bigr).
\label{eq:flatmoduli}
\end{equation}
当 $t$ 为实数时，方程变成椭圆型方程族；$t=0,\infty$ 是退化极限，非平凡瞬子只在极限附近出现。\src{B p.26}

\subsection{一族理论：$t\in\PP^1$ 与典范参数 $\Psi$}

由 \eqref{eq:Q-combo}，只差一个整体标度，两个数 $(u,v)$ 给出一个复参数
\begin{equation}
t=\frac{v}{u}\in\PP^1=\CC\cup\{\infty\}.
\label{eq:t-param}
\end{equation}
所以扭化不是给一个理论，而是给一族理论。族上还有两个「显然」的等价：整体缩放 $(u,v)\mapsto\lambda(u,v)$ 不改变 $t$；$SU(4)_R$ 的中心元素 $I$ 作用为 $I(t)=-t$，把 $t$ 与 $-t$ 等同（\src{B p.22}）。除这两条之外，四维一侧没有别的等价——真正的等价来自 S-对偶，因子 $(c\tau+d)/|c\tau+d|$ 作用在 $t$ 上，使 $t$ 不再只是标签。这正是后面能生出几何 Langlands 的地方。

物理上真正有意义的参数不是 $t$，而是 $t$ 与复耦合 $\tau=\theta/2\pi+4\pi\ii/g_{\mathrm{YM}}^2$ 的组合。原因是：扭化后作用量可以写成
\begin{equation}
I=\delta_t V+\frac{\ii\Psi}{4\pi}\int_M\Tr(F\wedge F),
\label{eq:action-Psi}
\end{equation}
第一项是超对称变分，第二项是拓扑项，其系数只可能依赖一个组合。配平量纲与 $\tau$ 依赖后得到
\begin{equation}
\Psi=\frac{\tau+\bar\tau}{2}+\frac{\tau-\bar\tau}{2}\cdot\frac{t-t^{-1}}{t+t^{-1}}
=\operatorname{Re}\tau+\ii\operatorname{Im}\tau\,\frac{t-t^{-1}}{t+t^{-1}}.
\label{eq:Psi}
\end{equation}
\src{B pp.35--36} 把 \eqref{eq:Psi} 右边通分，可以写成等价的有理形式
\begin{equation}
\Psi=\frac{\tau t^2+\bar\tau}{1+t^2}.
\label{eq:Psi-rational}
\end{equation}
这个形式在讨论特殊点时最方便。它的逆（\src{B p.36 eq (3.52)}）是
\begin{equation}
t^2=\frac{\bar\tau-\Psi}{\Psi-\tau},
\label{eq:t2-inverse}
\end{equation}
注意这里没有写错：分子分母都含 $\tau$ 的共轭差，而不是把 $\Psi-\tau$ 倒过来。我们核对过原文页图像；某些转录版本把共轭符号丢了，变成错误的倒数式，读者若按错误式子验算会得到相反的 $t$ 依赖。

由 \eqref{eq:Psi-rational} 立刻读出几个特殊点。$t=\pm\ii$ 时分母 $1+t^2=0$，而分子 $\tau t^2+\bar\tau=\bar\tau-\tau\neq0$（只要 $\operatorname{Im}\tau\neq0$），所以
\begin{equation}
t=\pm\ii\quad\Longrightarrow\quad \Psi=\infty,
\label{eq:Psi-inf}
\end{equation}
并且这个 $\infty$ 与 $\tau$ 的具体值无关。注意论证用的是有理形式：若用逆式 \eqref{eq:t2-inverse}，在 $t=\pm\ii$ 处会遇上 $0/0$，什么也读不出来。$t=\pm1$ 时 $(t-t^{-1})/(t+t^{-1})=0$，故 $\Psi=\operatorname{Re}\tau$；$|t|=1$ 时 $t=\ee^{\ii\theta}$，$(t-t^{-1})/(t+t^{-1})=\ii\tan\theta$，于是 $\Psi=\operatorname{Re}\tau-\operatorname{Im}\tau\tan\theta$ 为实数；弱耦合 $\operatorname{Im}\tau\to\infty$ 而固定 $\Psi,\operatorname{Re}\tau$ 时，由逆式得 $t^2\to1$，即 $t\to\pm1$。\src{B pp.36--37}

为什么说理论只依赖 $\Psi$，即为什么 $t$ 与 $\tau$ 的其它组合都不进入物理？原文给了两步论证（\src{B pp.36--37}）。第一步，作用量 \eqref{eq:action-Psi} 里只有拓扑项的系数含 $\tau$ 与 $t$，而所有可观测的 $\tau,t$ 依赖都必须通过这个系数，也就是通过 $\Psi$。第二步，对 $t\neq\pm\ii$ 定义
\begin{equation}
\delta'_t=\frac{\delta_t}{\sqrt{1+t^2}},
\label{eq:delta-prime}
\end{equation}
则 $(\delta'_t)^2=-\ii L_\sigma$ 与 $t$ 无关；再用小的场重定义把作用量整理成 Cartan 模型的标准形式，$t$ 的剩余依赖也被吸收掉。两步合起来：$(t,\tau)$ 只有通过 $\Psi$ 才有物理意义。顺带解释了 $t\sim-t$：$\Psi$ 只含 $t^2$，所以 $t$ 与 $-t$ 给出同一个 $\Psi$。

在 $\Psi$ 的语言里，S-对偶与 T-变换的作用极其简单：
\begin{equation}
S:\ \Psi\mapsto-\frac{1}{n_g\Psi},\qquad
T:\ \Psi\mapsto \Psi+1,
\label{eq:S-T}
\end{equation}
一般群元素 $\gamma=\begin{pmatrix}a&b\\c&d\end{pmatrix}\in SL(2,\ZZ)$ 给出 $\Psi\mapsto(a\Psi+b)/(c\Psi+d)$。\src{B p.37} 这里 $n_g$ 是\term{lacing number}{花边数}：单边（ADE）时 $n_g=1$，$B,C,F_4$ 时 $n_g=2$，$G_2$ 时 $n_g=3$。它来自 $\Ldual$ 的根系与原根系的对偶，$S$ 的作用因此带一个分数系数；$n_g=1$ 时退化为熟悉的 $\Psi\mapsto-1/\Psi$。$S$ 把 $\Psi=\infty$ 送到 $\Psi=0$，这一条第 4 讲要用。

参数空间的形状见图 \ref{fig:cp1}。
\begin{figure}[htbp]
\centering
\begin{tikzpicture}[scale=1.0]
  \draw[thick] (0,0) circle (1.8);
  \node[below] at (0,-1.98) {$t=0$};
  \node[above] at (0,1.98) {$t=\infty$};
  \node[right] at (1.88,0) {$t=1$};
  \node[left] at (-1.88,0) {$t=-1$};
  \fill (0,0.55) circle (2.4pt);
  \node[right] at (0.12,0.62) {$t=\ii$};
  \draw (0,-0.55) circle (2.4pt);
  \draw[thick] (-0.07,-0.62)--(0.07,-0.48);
  \draw[thick] (-0.07,-0.48)--(0.07,-0.62);
  \node[right] at (0.12,-0.68) {$t=-\ii$};
  \node[align=left,anchor=west,text width=5.7cm] at (2.3,1.15)
    {$\ii,-\ii$：$\Psi=\infty$，B-模型\\ 与 $\tau$ 无关};
  \node[align=left,anchor=west,text width=5.7cm] at (2.3,-1.25)
    {实轴 $\pm1$：$\Psi=\operatorname{Re}\tau$；\\ 其余实 $t$：$\Psi$ 为实数，A-模型};
\end{tikzpicture}
\caption{参数空间 $\PP^1_t$ 的示意图：圆是 $\PP^1$ 的轮廓，实轴画成轮廓上的四个标记点，$\pm\ii$ 画在圆内（分别用实心点与叉点表示）。$t$ 与 $-t$ 给出同一个 $\Psi$，$t=\pm\ii$ 对应 $\Psi=\infty$。}
\label{fig:cp1}
\end{figure}

最后是一个必须说在前面的警告。这套构造只在 \emph{Euclidean 号差}下成立。在 Lorentz 号差下，扭化要用一个紧致群 $\mathrm{Spin}(6)$ 去「吃掉」非紧的 $\mathrm{Spin}(1,3)$，找不到这样的单射；若改用 $\mathrm{Spin}(1,5)$，得到的哈密顿量不是厄米的，能量无下界。原文在 \src{B p.56 脚注 19} 把这一点讲得很直白：扭化破坏幺正性，只在 Euclidean 号差下有意义。这不是技术细节，而是整个几何 Langlands 故事必须在复几何（Riemann 面、复结构）里讲的根本原因；读原文时凡遇到「流形」「度量」都要默认其欧几里得。另外，\src{B p.16 脚注 5} 指出：当 $\theta=0$ 且 $\Psi=0,\infty$ 时，连不可定向的流形也可以允许。还要提醒：本文的推导是物理推导而非数学证明，且只处理非分歧情形；2024 年前后 Gaitsgory 等人的代数几何工作才给出几何 Langlands 的严格证明。

\begin{keybox}
扭化不是改变物理，而是换一个在弯曲流形上也存在的标量超荷：GL 扭化取 $\mathrm{Spin}(4)\subset\mathrm{Spin}(4)\times\mathrm{Spin}(6)$ 的对角子群，使六个标量重排为一个一形式与两个标量，并留下一族以 $t\in\PP^1$（等价地以 $\Psi$）参数化的拓扑场论。
\end{keybox}

\begin{follow}
\begin{itemize}
\item \textbf{B §3.1, pp.17--24}：三种扭化、$\kappa$ 的选取、$Q^2=0$ 的论证、场的重排与 $\delta_T\sigma=0$（重点读 pp.17--19 与 p.24）。
\item \textbf{B §3.2, pp.25--28}：扭化后的方程 \eqref{eq:bps}、$t=\pm\ii$ 与平坦 $G_{\mathbb C}$ 联络、实 $t$ 的椭圆族。
\item \textbf{B §3.5, pp.35--38}：$\Psi$ 的定义式、t--$\tau$ 合并、$t=\pm\ii$ 的 $\Psi=\infty$、$S$ 与 $T$ 的作用。
\end{itemize}
\end{follow}

\begin{checkbox}
\begin{itemize}
\item 为什么平坦空间上的 16 个超荷不能直接搬到一般弯曲流形上？扭化究竟改变了什么、又没有改变什么？
\item 写出 $\mathbf 4$ 与 $\mathbf 6_v$ 在 $\mathrm{Spin}'(4)$ 下的分解，并解释 $\sigma$ 为何满足 $\delta_T\sigma=0$。
\item 用 \eqref{eq:Psi-rational} 验证 $t=\pm\ii$ 时 $\Psi=\infty$ 且与 $\tau$ 无关；说明为什么不能改用逆式 \eqref{eq:t2-inverse} 来验证。
\item 若两个不变超荷的 $K$ 荷不同，$(u,v)$ 的参数空间会变成什么？为什么？
\item 说明为什么可观测物理只依赖 $\Psi$，以及 $t$ 与 $-t$ 为何给出同一个理论。
\item 为什么这套扭化只在 Euclidean 号差下成立？Lorentz 号差下具体卡在哪一步？
\end{itemize}
\end{checkbox}
